\documentclass[11pt,a4paper]{article}

\usepackage[utf8]{inputenc}
\usepackage[T1]{fontenc}
\usepackage[french]{babel}
\usepackage[margin=2.3cm]{geometry}
\usepackage{lmodern}
\usepackage{microtype}
\usepackage{amsmath,amssymb,amsthm,mathtools}
\usepackage{bm}
\usepackage{graphicx}
\usepackage{booktabs}
\usepackage{caption}
\captionsetup{font=small,labelfont=bf}
\usepackage{subcaption}
\usepackage{float}
\usepackage{enumitem}
\usepackage{xcolor}
\usepackage{hyperref}
\usepackage{algorithm}
\usepackage{algpseudocode}
\usepackage{titlesec}

\hypersetup{
  colorlinks=true,
  linkcolor=black,
  citecolor=black,
  urlcolor=blue!50!black,
}

\setlength{\parindent}{1.2em}
\setlength{\parskip}{0.25em}
\renewcommand{\baselinestretch}{1.08}

\titleformat{\section}{\Large\bfseries}{\thesection}{0.7em}{}
\titleformat{\subsection}{\large\bfseries}{\thesubsection}{0.6em}{}
\titlespacing*{\section}{0pt}{1.4em}{0.5em}
\titlespacing*{\subsection}{0pt}{1em}{0.3em}

\newcommand{\E}{\mathbb{E}}
\newcommand{\Pb}{\mathbb{P}}
\newcommand{\IG}{\mathrm{IG}}

\title{\vspace{-1.5cm}\textbf{Modélisation stochastique d'un carnet d'ordres}\\[0.3em]
       \large MODAL --- Carnet d'ordres SNA}
\author{}
\date{}

\begin{document}
\maketitle
\vspace{-2cm}

\section*{Présentation}

Le projet vise à comprendre, par la simulation, comment les volumes au meilleur Bid
($q_b$) et au meilleur Ask ($q_a$) d'un carnet d'ordres évoluent au cours du
temps, et en particulier comment se distribue le premier instant où l'un des
deux côtés s'épuise --- événement microstructurel qui déclenche un mouvement
de prix. Notre démarche suit une progression de modèles, partant du plus simple
(processus de Poisson indépendants pour chaque côté) pour aller vers des
spécifications plus réalistes intégrant l'auto-excitation par processus de
Hawkes, puis le couplage entre les deux côtés du carnet, et enfin un troisième
processus représentant le niveau adjacent.

Pour chaque modèle on s'intéresse principalement à deux quantités~: le temps
moyen d'épuisement $\E[\min(\tau_a, \tau_b)]$ et la probabilité asymétrique
$\Pb(\tau_a < \tau_b)$. Ce rapport reprend, dans le même ordre que le notebook
\texttt{carnet\_d\_ordres.ipynb}, les méthodes mises en place, les choix
numériques, et les résultats obtenus, en discutant à chaque fois leur
crédibilité (intervalles de confiance, comparaison à des prédictions
théoriques quand elles existent).

\section{Cadre paramétrique}
\label{sec:params}

Avant d'entrer dans les modèles, il est utile de fixer les valeurs numériques
qui seront utilisées partout. Le tableau~\ref{tab:params} les rassemble. La
plupart de ces choix répondent à un objectif simple~: rester dans un régime où
l'atteinte de zéro est presque sûre (ce qui impose un drift négatif), tout en
gardant des temps de simulation raisonnables.

\begin{table}[H]
\centering
\small
\begin{tabular}{@{}lcp{0.58\linewidth}@{}}
\toprule
Paramètre & Valeur & Rôle / justification \\
\midrule
$\lambda^+$ & $1{,}2$ & Intensité d'insertion ; choisie inférieure à $\lambda^-$ pour garantir un drift négatif. \\
$\lambda^-$ & $2{,}0$ & Intensité d'annulation ; donne $\mu = -0{,}8$ et $\sigma^2 = 3{,}2$. \\
$T$ & $100$ & Horizon temporel des trajectoires de référence. \\
$q_0$ & $10$ & Volume initial : assez grand pour être dans le régime fluide, assez petit pour que $\tau$ reste observable. \\
$\mu^-$ & $2{,}0$ & Intensité de fond du Hawkes, prise égale à $\lambda^-$ pour comparer Poisson et Hawkes au même régime de fond. \\
$\alpha, \beta$ & $0{,}5\,,\,1{,}0$ & Gain et taux de relaxation Hawkes ; $\alpha/\beta = 1/2$ assure la stabilité. \\
$a, b$ & $0{,}5\,,\,1{,}0$ & Mêmes échelles pour l'excitation de $q_2$ par les décréments d'ask. \\
$dt$ & $0{,}01$ & Pas du schéma d'Euler-Maruyama pour le mouvement brownien. \\
$M$ (espérances) & $10^3$ à $5\cdot10^3$ & $\hat\sigma/\sqrt{M}$ de l'ordre de $1\,\%$ de $\hat T$. \\
$M$ (surfaces 3D) & $200$ à $500$ & Compromis temps de calcul / précision sur grille $10 \times 10$. \\
\bottomrule
\end{tabular}
\caption{Paramètres globaux. Les valeurs des intensités sont volontairement
modestes pour conserver des trajectoires lisibles.}
\label{tab:params}
\end{table}

Trois conséquences théoriques immédiates serviront ensuite de référence pour
juger les estimateurs Monte-Carlo~:

\[
\E[\tau]_{\text{th}} = \frac{q_0}{|\mu|} = 12{,}5,
\qquad
\mathrm{Var}(\tau)_{\text{th}} = \frac{q_0\,\sigma^2}{|\mu|^3} = 62{,}5,
\qquad
\bar\lambda^-_{\infty} = \frac{\mu^-\beta - \lambda^+\alpha}{\beta - \alpha} = 2{,}8.
\]

\section{Modèle Poisson de base}

On commence par le modèle le plus élémentaire. Chaque côté du carnet est décrit
par
\[
q(t) = q(0) + N^+_t - N^-_t,
\]
où $N^+$ et $N^-$ sont deux processus de Poisson indépendants d'intensités
$\lambda^+$ et $\lambda^-$, et où le processus est absorbé en zéro. La
simulation est exacte~: les inter-événements suivent une loi
$\mathrm{Exp}(\lambda^+ + \lambda^-)$, et chaque événement est un saut $+1$ avec
probabilité $\lambda^+ / (\lambda^+ + \lambda^-)$, sinon $-1$. Pour étudier les
deux côtés simultanément, on lance deux instances indépendantes.

La figure~\ref{fig:traj} montre une réalisation typique. Le drift négatif est
visible (la trajectoire dérive vers le bas) mais des excursions vers le haut
restent fréquentes, ce qui rend le temps d'atteinte de zéro lui-même très
variable.

\begin{figure}[H]
\centering
\includegraphics[width=0.85\linewidth]{figures/trajectories_subplots.png}
\caption{Réalisation de $q_a$ (haut) et $q_b$ (bas) jusqu'à atteinte de zéro,
avec $\lambda^+ = 1{,}2$, $\lambda^- = 2$, $q_0 = 10$. Ici l'ask se vide
légèrement avant le bid.}
\label{fig:traj}
\end{figure}

\section{Estimation du temps moyen d'atteinte}

Pour estimer $\E[\tau]$, on tire $M$ trajectoires indépendantes et on prend la
moyenne empirique~:
\[
\widehat T_M = \frac{1}{M} \sum_{i=1}^M \tau^{(i)}, \qquad
\text{IC}_{1-\alpha} = \widehat T_M \pm z_{1-\alpha/2}\, \frac{\widehat\sigma_M}{\sqrt{M}}.
\]
Le théorème central limite garantit la validité asymptotique de l'intervalle de
confiance, à condition que les $\tau^{(i)}$ aient une variance finie, ce qui
est le cas dès que le drift est strictement négatif.

Sur $M = 2000$ trajectoires, on obtient
\[
\widehat T_{\text{1 proc.}} = 12{,}25 \quad \text{IC}_{95\%} = [11{,}90 \,;\, 12{,}60],
\qquad
\widehat T_{\min} = 8{,}53 \quad \text{IC}_{95\%} = [8{,}33 \,;\, 8{,}72].
\]
Le premier est l'espérance du temps d'atteinte d'un seul côté, le second celle
du minimum entre l'ask et le bid. La référence brownienne $q_0/|\mu| = 12{,}5$
tombe à l'intérieur de l'intervalle de confiance --- c'est rassurant pour la
suite, et cela confirme que l'approximation diffusive (qu'on validera plus
loin) sait déjà capturer la moyenne. Le passage au minimum réduit cette moyenne
d'un facteur $\approx 1{,}45$, comme on s'y attend pour le minimum de deux
variables IG indépendantes de mêmes paramètres.

\section{Trajectoire conjointe dans le plan $(q_b, q_a)$}

Plutôt que de regarder $q_a$ et $q_b$ séparément en fonction du temps, on peut
représenter leur trajectoire conjointe dans le plan, ce qui permet de
visualiser quels couples d'états sont visités. Les deux processus n'ayant pas
les mêmes temps d'événement, on synchronise les deux échelles par recherche
binaire (\texttt{np.searchsorted}) sur les temps fusionnés ; le coût est
$O(n\log n)$. La transparence des points encode la fréquence d'occupation~: les
états récurrents apparaissent opaques, les transitoires plus pâles.

\begin{figure}[H]
\centering
\includegraphics[width=0.62\linewidth]{figures/trajectory_plane.png}
\caption{Trajectoire dans le plan $(q_b, q_a)$ jusqu'à ce qu'un côté atteigne
zéro. Le point rouge marque le départ, le point vert l'arrêt.}
\label{fig:plane}
\end{figure}

Pour les valeurs choisies ($\lambda^- > \lambda^+$), la trajectoire dérive
assez nettement vers l'origine. Une utilisation plus quantitative de cette
représentation pourrait passer par l'estimation d'un noyau de densité
2D~; on ne l'a pas faite ici, le caractère illustratif du graphique suffisant
à ce stade.

\section{Approximation diffusive}

Lorsque $\lambda^\pm$ sont grandes (régime fluide), un théorème central limite
fonctionnel pour les processus de Poisson centrés-renormalisés montre que $q$
peut être approché par un mouvement brownien~:
\[
\mathrm{d}Q_t = (\lambda^+ - \lambda^-)\,\mathrm{d}t
              + \sqrt{\lambda^+ + \lambda^-}\,\mathrm{d}W_t,
\quad Q_0 = q_0.
\]
À horizon $T$ fixé (en l'absence d'absorption), il en découle
\[
q(T) \;\sim\; \mathcal{N}\!\bigl(q_0 + \mu T,\; \sigma^2 T\bigr).
\]

Sur $5000$ trajectoires libres jusqu'à $T = 100$, la moyenne empirique vaut
$-70{,}56$ contre $-70$ théorique, et l'écart-type empirique $17{,}74$ contre
$\sigma\sqrt{T} \approx 17{,}89$. Les deux moments sont donc reproduits à mieux
de $1\,\%$. La figure~\ref{fig:gauss} superpose l'histogramme à la densité
gaussienne théorique~; l'accord visuel est sans équivoque.

\begin{figure}[H]
\centering
\includegraphics[width=0.62\linewidth]{figures/gaussian_approx.png}
\caption{Distribution empirique de $q(T)$ contre $\mathcal{N}(-70, 320)$. Sans
absorption, la marche aléatoire entière est très bien approchée par sa limite
brownienne dès $T = 100$.}
\label{fig:gauss}
\end{figure}

\section{Loi du temps d'atteinte}

Une conséquence classique de l'approximation diffusive est que le premier
temps d'atteinte de zéro d'un mouvement brownien $X_t = q_0 + \mu t + \sigma W_t$
de drift négatif suit une loi inverse-gaussienne~:
\[
\tau \sim \IG(\mu_{\IG}, \lambda_{\IG}),
\qquad
\mu_{\IG} = \frac{q_0}{|\mu|},
\qquad
\lambda_{\IG} = \frac{q_0^2}{\sigma^2}.
\]

Pour vérifier que cette prédiction reste valable pour la marche aléatoire
entière qui nous intéresse, on simule $N = 5000$ trajectoires Poisson, on
calcule le temps d'arrêt de chacune, puis on compare la distribution empirique
à la densité IG théorique. Deux indicateurs sont rapportés~: l'accord des deux
premiers moments (table~\ref{tab:igval}) et un test de Kolmogorov-Smirnov.

\begin{table}[H]
\centering
\small
\begin{tabular}{@{}lccc@{}}
\toprule
& Théorique (BM) & Empirique (Poisson) & Écart relatif \\
\midrule
Moyenne $\E[\tau]$            & $12{,}500$ & $12{,}537$ & $+0{,}30\,\%$ \\
Variance $\mathrm{Var}(\tau)$ & $62{,}500$ & $66{,}731$ & $+6{,}77\,\%$ \\
\midrule
Test KS & \multicolumn{3}{c}{$D = 0{,}013$, \quad $p\text{-value} = 0{,}369$} \\
\bottomrule
\end{tabular}
\caption{Adéquation de la loi inverse-gaussienne aux temps d'arrêt simulés.}
\label{tab:igval}
\end{table}

\begin{figure}[H]
\centering
\includegraphics[width=0.62\linewidth]{figures/ig_distribution.png}
\caption{Histogramme des temps d'arrêt simulés et densité IG théorique. Le test
de Kolmogorov-Smirnov ne rejette pas l'hypothèse d'adéquation.}
\label{fig:igval}
\end{figure}

L'adéquation est excellente sur la moyenne ($+0{,}3\,\%$), un peu moins bonne
sur la variance ($+6{,}8\,\%$). Cet écart sur la variance n'est pas surprenant~:
la marche entière est plus bruitée localement qu'une diffusion continue, et
$q_0 = 10$ n'est pas grand au sens asymptotique. Le test KS, qui ne dépend pas
spécifiquement des moments, donne $p = 0{,}37$~: on ne peut donc pas distinguer
la distribution empirique de la loi théorique au seuil habituel de $5\,\%$.
L'approximation IG est validée.

\section{Probabilité que l'ask atteigne zéro avant le bid}

Lorsque les deux côtés sont indépendants et de mêmes intensités, la quantité
$p(q_0^a, q_0^b) = \Pb(\tau_a < \tau_b)$ est, dans la limite continue, donnée
par
\[
p(q_0^a, q_0^b) = \int_0^\infty f_{\tau_a}(t)\,\bigl(1 - F_{\tau_b}(t)\bigr)\,\mathrm{d}t,
\]
avec $f$ et $F$ les densité et fonction de répartition inverse-gaussiennes.
Cette intégrale n'a pas de forme close mais se calcule numériquement très
rapidement par quadrature.

On compare cette prédiction à la surface obtenue par simulation Monte-Carlo
sur la grille $(q_0^a, q_0^b) \in \{1,\dots,10\}^2$ avec $M = 500$
trajectoires par grain. Le résultat est en figure~\ref{fig:ruinval}.

\begin{figure}[H]
\centering
\includegraphics[width=\linewidth]{figures/ruin_validation.png}
\caption{Surface MC (gauche), prédiction par quadrature sur les IG (centre),
écart entre les deux (droite). L'erreur reste dans la marge attendue du bruit
Monte-Carlo.}
\label{fig:ruinval}
\end{figure}

L'erreur absolue maximale sur la grille vaut $0{,}050$ et l'erreur moyenne
$0{,}015$. Pour $M = 500$, le bruit Monte-Carlo théorique sur une probabilité
est borné par $\sqrt{p(1-p)/M} \le 0{,}022$. L'écart MC vs prédiction
brownienne est donc principalement explicable par ce bruit, et il n'y a pas
de biais systématique détectable. Sur la diagonale ($q_0^a = q_0^b$), la
moyenne empirique de $p$ vaut $0{,}497$, conformément au $1/2$ qu'imposerait
une symétrie parfaite.

\section{Surface 3D du temps moyen d'atteinte}

Dans le même esprit que la surface précédente, on regarde maintenant la
surface du temps moyen $T(q_0^a, q_0^b) = \E[\min(\tau_a, \tau_b)]$, avec son
intervalle de confiance à 95\,\% (calculé par TCL pour chaque grain). La
figure~\ref{fig:poisson3d} affiche, côte à côte, la surface des temps moyens
(à gauche, avec les bornes d'IC en rouge et noir) et celle des probabilités
déjà discutée plus haut (à droite). Cette double représentation permet de
visualiser le \emph{quand} et le \emph{lequel}~: à l'angle inférieur gauche
l'épuisement est rapide mais incertain ; quand on s'approche de la diagonale
et de la grande taille $(10,10)$, le temps moyen approche $8{,}3$ unités et
$p \approx 1/2$.

\begin{figure}[H]
\centering
\includegraphics[width=\linewidth]{figures/time_proba_3d_poisson.png}
\caption{Modèle Poisson indépendant. Gauche~: temps moyen $\E[\min(\tau_a,
\tau_b)]$ avec IC 95\,\% (wireframes rouge et noire). Droite~: probabilité que
l'ask atteigne zéro en premier.}
\label{fig:poisson3d}
\end{figure}

L'intervalle de confiance reste serré sur l'ensemble de la grille (typiquement
$\pm 2\,\%$ de la valeur moyenne pour $M = 300$), ce qui rend la surface
suffisamment lisse pour des extractions futures (par exemple, pour calibrer un
critère d'agressivité d'exécution en fonction des volumes courants).

\section{Processus de Hawkes individuel}

On enrichit le modèle en remplaçant l'intensité constante $\lambda^-$ par une
intensité de Hawkes auto-excitée~:
\[
\lambda^-(t) = \mu^- + \sum_{T_i < t} \alpha\, e^{-\beta(t - T_i)}.
\]
Concrètement, chaque cancellation augmente instantanément $\lambda^-$ de
$\alpha$, qui décroît ensuite exponentiellement à un taux $\beta$ vers la
valeur de fond $\mu^-$. On modélise par ailleurs l'effet inverse pour les
insertions (qui font baisser $\lambda^-$ de $\alpha$).

En prenant l'espérance de l'EDS sur $\lambda^-$ et en cherchant un point fixe
stationnaire, on obtient
\begin{equation}
\bar\lambda^-_\infty = \frac{\mu^-\beta - \lambda^+\alpha}{\beta - \alpha},
\label{eq:lambda_stat}
\end{equation}
ce qui n'est fini que si $\alpha < \beta$. Avec nos valeurs, $\bar\lambda^-_\infty
= 2{,}8$.

\paragraph{Simulation par thinning d'Ogata.}
L'intensité $\lambda^-(t)$ étant non bornée \emph{a priori} mais bornée par
morceaux entre deux événements, on simule par l'algorithme classique de
thinning. Entre deux événements on majore $\lambda^-(t)$ par $\bar\lambda$, on
échantillonne un candidat exponentiel de paramètre $\bar\lambda + \lambda^+$,
et on accepte ou rejette selon la valeur effective de l'intensité au candidat.

\begin{algorithm}[H]
\caption{Thinning d'Ogata pour le Hawkes individuel}
\begin{algorithmic}[1]
\State $t \gets 0$,~ $\lambda^- \gets \mu^-$
\While{condition d'arrêt non atteinte}
  \State $\bar\lambda \gets \max(\mu^-, \lambda^-)$,~ $\Lambda \gets \bar\lambda + \lambda^+$
  \State Tirer $\Delta t \sim \mathcal{E}(\Lambda)$,~ $t' \gets t + \Delta t$
  \State $\lambda^-_{\mathrm{pre}} \gets \mu^- + (\lambda^- - \mu^-)e^{-\beta(t' - t_{\text{last}})}$
  \State Tirer $U \sim \mathcal{U}([0,1])$
  \If{$U < \lambda^+/\Lambda$} \Comment{insertion}
    \State $q \gets q + 1$,~ $\lambda^- \gets \lambda^-_{\mathrm{pre}} - \alpha$
  \ElsIf{$U < (\max(\lambda^-_{\mathrm{pre}}, 0) + \lambda^+)/\Lambda$} \Comment{annulation}
    \State $q \gets q - 1$,~ $\lambda^- \gets \lambda^-_{\mathrm{pre}} + \alpha$
  \Else
    \State $\lambda^- \gets \lambda^-_{\mathrm{pre}}$ \Comment{rejet, on relaxe seulement}
  \EndIf
  \State $t \gets t'$
\EndWhile
\end{algorithmic}
\end{algorithm}

La figure~\ref{fig:hawkes_traj} montre une trajectoire typique de $q$ et
l'évolution de $\lambda^-(t)$ correspondante, où l'on voit nettement les pics
d'intensité après chaque cancellation et la relaxation exponentielle entre
deux événements.

\begin{figure}[H]
\centering
\begin{subfigure}[t]{0.49\linewidth}
  \centering
  \includegraphics[width=\linewidth]{figures/hawkes_trajectory.png}
  \caption{$q(t)$ sur $[0, T]$.}
\end{subfigure}\hfill
\begin{subfigure}[t]{0.49\linewidth}
  \centering
  \includegraphics[width=\linewidth]{figures/hawkes_lambda.png}
  \caption{$\lambda^-(t)$ correspondante.}
\end{subfigure}
\caption{Trajectoire Hawkes individuelle et son intensité $\lambda^-(t)$.}
\label{fig:hawkes_traj}
\end{figure}

\paragraph{Validation de la moyenne stationnaire.}
On vérifie numériquement \eqref{eq:lambda_stat} en simulant cent trajectoires
sur un horizon long ($T = 2000$), en jetant un \emph{burn-in} de $200$ unités,
puis en calculant la moyenne temporelle de $\lambda^-(t)$ sur la fenêtre
post-burn-in. Le résultat~: la moyenne empirique vaut $2{,}97 \pm 0{,}01$ (IC
95\,\%) contre $2{,}80$ théorique, soit un écart relatif de $+6\,\%$. Cet
écart, statistiquement significatif, provient du \emph{clipping} de
l'intensité dans l'algorithme~: quand la relaxation entraîne $\lambda^-_{\mathrm
{pre}}$ en territoire négatif, on annule simplement les événements de
cancellation au lieu de leur attribuer une probabilité négative, ce qui
introduit un biais positif sur la moyenne stationnaire. La distribution de
$\lambda^-(T)$ obtenue (figure~\ref{fig:hawkes_lambda_T}) est unimodale et
centrée légèrement au-dessus de la valeur théorique, comme prévu.

\begin{figure}[H]
\centering
\includegraphics[width=0.6\linewidth]{figures/hawkes_lambda_T.png}
\caption{Distribution empirique de $\lambda^-(T)$ pour $T = 1000$ et la valeur
théorique $\bar\lambda^-_\infty = 2{,}8$ (trait rouge). Le mode empirique est
légèrement au-dessus, conformément à l'analyse du biais.}
\label{fig:hawkes_lambda_T}
\end{figure}

La distribution du temps d'arrêt sous Hawkes
(figure~\ref{fig:hawkes_stop}) confirme par ailleurs le comportement
attendu~: par rapport au cas Poisson, l'auto-excitation génère beaucoup de
trajectoires courtes (avalanches de cancellations) et étire la queue droite
(rare cas où l'excitation ne se déclenche pas).

\begin{figure}[H]
\centering
\includegraphics[width=0.6\linewidth]{figures/hawkes_stopping_dist.png}
\caption{Distribution empirique du temps d'arrêt sous Hawkes individuel
($\mu^- = 2$, $\alpha = 0{,}5$, $\beta = 1$).}
\label{fig:hawkes_stop}
\end{figure}

\section{Hawkes couplé bid/ask}

L'étape suivante consiste à coupler les deux côtés du carnet. On suit
conjointement $(q_a, q_b)$ avec deux intensités $(\lambda^-_a, \lambda^-_b)$
auto- et croise-excitées~: chaque événement sur l'ask diminue $\lambda^-_a$ et
augmente $\lambda^-_b$ (et symétriquement). Cela modélise la propagation d'un
déséquilibre~: lorsqu'un côté du carnet se vide, la pression de marché tend à
s'accumuler de l'autre côté.

L'algorithme de thinning d'Ogata est adapté avec une borne globale
$\bar\lambda_a + \bar\lambda_b + \lambda^+_a + \lambda^+_b$ et une mise à jour
conjointe des deux intensités à chaque événement effectif. La
figure~\ref{fig:coupled} présente une réalisation et la distribution
correspondante du temps d'arrêt.

\begin{figure}[H]
\centering
\begin{subfigure}[t]{0.49\linewidth}
  \centering
  \includegraphics[width=\linewidth]{figures/coupled_hawkes_traj.png}
  \caption{Trajectoires couplées $q_a$ et $q_b$.}
\end{subfigure}\hfill
\begin{subfigure}[t]{0.49\linewidth}
  \centering
  \includegraphics[width=\linewidth]{figures/coupled_hawkes_stopping.png}
  \caption{Distribution empirique de $\min(\tau_a, \tau_b)$.}
\end{subfigure}
\caption{Modèle de Hawkes couplé bid/ask, $\mu^-_a = \mu^-_b = 2$.}
\label{fig:coupled}
\end{figure}

\paragraph{Surfaces 3D pour le Hawkes couplé.}
On reproduit pour ce modèle les deux surfaces 3D précédentes
(figure~\ref{fig:hawkes3d}). Contrairement au cas Poisson indépendant,
il n'existe pas de prédiction analytique simple à comparer~: le couplage entre
$\lambda^-_a$ et $\lambda^-_b$ rend le système non-Markovien dans
$(q_a, q_b)$ seul (il faut aussi suivre l'état des intensités). Le Monte-Carlo
reste néanmoins le bon outil.

\begin{figure}[H]
\centering
\includegraphics[width=\linewidth]{figures/hawkes_coupled_3d.png}
\caption{Hawkes couplé bid/ask. Gauche~: temps moyen avec IC 95\,\%. Droite~:
probabilité que l'ask atteigne zéro en premier.}
\label{fig:hawkes3d}
\end{figure}

Au coin $(q_0^a, q_0^b) = (10, 10)$, on mesure $\E[\min(\tau_a, \tau_b)] =
5{,}83$ contre $8{,}32$ pour le Poisson équivalent, soit un facteur
$\approx 0{,}70$. L'auto-excitation accélère donc nettement l'épuisement, et la
moyenne diagonale de la probabilité reste $\approx 0{,}52$ (proche de $1/2$ par
symétrie des paramètres).

\section{Comparaison des quatre régimes}
\label{sec:cmp}

À ce stade on dispose de quatre modèles complets~: Poisson individuel, premier
de deux Poisson indépendants, Hawkes individuel, et Hawkes couplé. La
table~\ref{tab:cmp} et la figure~\ref{fig:cmp} les comparent sur le même jeu de
paramètres ($q_0 = 10$, $1000$ trajectoires sauf le Hawkes couplé limité à
$300$).

\begin{table}[H]
\centering
\small
\begin{tabular}{@{}lcccc@{}}
\toprule
Modèle & Moyenne & Écart-type & Médiane & IQR \\
\midrule
Poisson individuel             & $12{,}42$ & $7{,}41$ & $10{,}62$ & $[7{,}28\,;\,15{,}30]$ \\
Premier de deux Poisson indép. & $\phantom{0}8{,}56$ & $4{,}25$ & $\phantom{0}7{,}67$ & $[5{,}60\,;\,10{,}52]$ \\
Hawkes individuel              & $\phantom{0}8{,}80$ & $6{,}83$ & $\phantom{0}6{,}89$ & $[4{,}18\,;\,11{,}17]$ \\
Hawkes couplé bid/ask          & $\phantom{0}5{,}91$ & $3{,}55$ & $\phantom{0}5{,}11$ & $[3{,}36\,;\,7{,}42]$ \\
\bottomrule
\end{tabular}
\caption{Comparaison synthétique des distributions de temps d'arrêt.}
\label{tab:cmp}
\end{table}

\begin{figure}[H]
\centering
\includegraphics[width=0.85\linewidth]{figures/comparison.png}
\caption{Densités empiriques du temps d'arrêt pour les quatre régimes.}
\label{fig:cmp}
\end{figure}

Trois enseignements ressortent. D'abord, le passage au minimum de deux Poisson
indépendants divise la moyenne par environ $1{,}45$ et la variance par $3$~:
c'est un effet purement combinatoire, sans rien à voir avec la dynamique. La
distribution se concentre par construction. Ensuite, l'auto-excitation à un
seul processus (Hawkes individuel) fait baisser la médiane plus que la
moyenne~: la plupart des trajectoires meurent rapidement (avalanches), avec
une minorité qui résiste longtemps quand l'excitation ne se déclenche pas.
Enfin, le couplage bid/ask combine ces deux effets et amplifie l'asymétrie~:
chaque événement sur un côté excite l'autre côté en sens inverse, et l'on
obtient au final la distribution la plus concentrée et la moyenne la plus
courte des quatre.

\section{Processus auxiliaire $q_2$}

Une question naturelle pour la microstructure est de savoir ce qu'on observe
au niveau adjacent du carnet au moment précis où le meilleur niveau d'un côté
s'épuise. C'est l'objet du processus auxiliaire $q_2$, simulé en parallèle des
deux files Hawkes couplées. Dans une première version, $q_2$ est lui-même un
Poisson constant de paramètres $(\lambda^+_2, \lambda^-_2)$. On regarde la
distribution conditionnelle de $q_2(\tau)$ selon que l'événement déclenchant
$\tau$ soit l'épuisement de l'ask ($\tau = \tau_a$) ou celui du bid ($\tau =
\tau_b$).

La figure~\ref{fig:q2_cond} montre le résultat dans le cas \emph{asymétrique}
($\mu^-_a = 2$, $\mu^-_b = 1{,}2$), c'est-à-dire avec un ask plus fragile que
le bid. La distribution conditionnelle à $\tau = \tau_a$ se trouve nettement
décalée vers le haut~: lorsque l'ask cède en premier, c'est typiquement après
un délai court, et $q_2$ a eu peu de temps pour dériver vers le bas. À
l'inverse, conditionner sur $\tau = \tau_b$ sélectionne des trajectoires plus
longues (le côté faible n'a pas cédé), pendant lesquelles $q_2$ a eu le temps
de baisser.

\begin{figure}[H]
\centering
\includegraphics[width=\linewidth]{figures/q2_asymmetric.png}
\caption{Distribution conditionnelle de $q_2(\tau)$ selon le côté qui s'épuise
en premier, dans le cas asymétrique $\mu^-_a = 2$, $\mu^-_b = 1{,}2$. La
sélection induit un effet visible sur la distribution.}
\label{fig:q2_cond}
\end{figure}

Dans le cas symétrique ($\mu^-_a = \mu^-_b = 2$), les deux distributions sont
en revanche presque identiques (figure~\ref{fig:q2_sym}, à gauche), ce qui est
attendu~: par symétrie, conditionner sur le côté gagnant ne change rien à la
loi marginale d'un processus indépendant.

\section{Intensité $\lambda^+_2$ excitée par les décréments d'ask}

On raffine maintenant le modèle de $q_2$ en couplant son intensité d'insertion
au flux des événements décrémentant l'ask~:
\[
\lambda^+_2(t) = \mu^+_2 + \sum_{T_i^{\downarrow,a} < t} a\, e^{-b(t - T_i^{\downarrow,a})}.
\]
L'idée est qu'une consommation rapide d'un niveau d'ask augmente la probabilité
d'apparition de nouveaux ordres au niveau adjacent. La simulation reprend le
même thinning, en majorant $\lambda^+_2$ par $\mu^+_2 + a \times (\text{nombre
total de baisses d'ask sur le run de référence})$.

La figure~\ref{fig:q2_sym} compare, dans le cas symétrique, la distribution
conditionnelle obtenue sous le modèle Poisson constant (à gauche, déjà vue
plus haut) et sous le modèle excité (à droite). Le déplacement vers la droite
des distributions excitées est cohérent avec l'intuition~: une cancellation
côté ask provoque, en moyenne, plus d'insertions sur $q_2$ qu'auparavant.

\begin{figure}[H]
\centering
\begin{subfigure}[t]{0.49\linewidth}
  \centering
  \includegraphics[width=\linewidth]{figures/q2_symmetric_poisson.png}
  \caption{$q_2$ Poisson constant (référence).}
\end{subfigure}\hfill
\begin{subfigure}[t]{0.49\linewidth}
  \centering
  \includegraphics[width=\linewidth]{figures/q2_symmetric_excited.png}
  \caption{$q_2$ avec $\lambda^+_2$ excitée par les décréments d'ask.}
\end{subfigure}
\caption{Distribution conditionnelle de $q_2(\tau)$ dans le cas symétrique
($\mu^-_a = \mu^-_b = 2$). À gauche, modèle indépendant ; à droite, modèle
couplé via les décréments d'ask.}
\label{fig:q2_sym}
\end{figure}

L'effet est visible mais reste, à ce stade, qualitatif. Une suite naturelle
serait d'ajouter un test statistique formel (par exemple un test de
Mann-Whitney sur les deux distributions) pour quantifier la significativité
de l'écart entre la version Poisson et la version excitée, ainsi qu'une étude
de sensibilité aux paramètres $a$ et $b$.

\section{Conclusion et pistes}

Sur l'ensemble des modèles considérés, l'estimation Monte-Carlo donne des
résultats cohérents et, là où une référence théorique existe, en bon accord
avec celle-ci~: la moyenne et la variance du temps d'atteinte sont reproduites
à mieux de $1\,\%$ et $7\,\%$ respectivement par rapport à la prédiction IG,
le test de Kolmogorov-Smirnov ne rejette pas l'hypothèse d'adéquation à la
loi inverse-gaussienne, et la surface MC de $\Pb(\tau_a < \tau_b)$ est
indistinguable de la prédiction brownienne au niveau du bruit. Le seul point
de désaccord notable concerne la moyenne stationnaire de $\lambda^-$ sous
Hawkes, qui présente un biais positif d'environ $6\,\%$ imputable au
\emph{clipping} d'intensité de l'algorithme~; ce biais a été identifié et son
origine est claire.

Le passage du Poisson au Hawkes, puis du Hawkes individuel au Hawkes couplé,
divise environ par deux le temps moyen d'épuisement à $q_0 = 10$ (de $12{,}4$
à $5{,}9$). Les distributions correspondantes sont aussi de plus en plus
concentrées. Cela illustre quantitativement à quel point l'auto-excitation et
le couplage entre les deux côtés du carnet modifient le risque
d'épuisement~: ignorer l'un ou l'autre conduit à \emph{sous-estimer}
significativement la probabilité d'un mouvement de prix à court terme.

Plusieurs pistes restent ouvertes. La plus immédiate est l'étude systématique
du cas asymétrique ($\mu^-_a \neq \mu^-_b$), pour lequel la surface
$\Pb(\tau_a < \tau_b)$ devient un véritable outil de quantification de
l'imbalance. Le couplage du processus auxiliaire $q_2$ avec les événements
d'ask mériterait également un test statistique en règle, ainsi qu'une étude
de sensibilité aux paramètres d'excitation $(a, b)$. Enfin, la calibration
sur données réelles (intensités $\lambda^\pm$, paramètres Hawkes) permettrait
de passer d'une étude qualitative à un véritable outil prédictif d'aide à
l'exécution.

\end{document}
